% !PRINTED 171
\weberpage{171}
Endlich können wir auch noch nach gebrochenen Invarianten fragen. Ist
\[
\frac{F(x)}{F_1(x)}
\]
eine solche gebrochene Invariante, so nehmen wir zunächst an, die beiden ganzen rationalen Functionen $F(x), F_1(x)$ von den $n$ Variablen $x$ seien von gemeinschaftlichen Theilern befreit (Bd.\,I, \S.\,51).

Beschränken wir uns fürs Erste auf absolute Invarianten, und ist demnach
\[
\frac{F(x)}{F_1(x)} = \frac{F[A(x)]}{F_1[A(x)]},
\]
so haben, da die $x$ hier als unabhängige Variable angesehen werden, weder rechts noch links Zähler und Nenner einen gemeinschaftlichen Theiler, und es folgt:
\[
F[A(x)] = \alpha F(x), \quad F_1[A(x)] = \alpha F_1(x),
\]
worin $\alpha$ ein constanter Factor ist, der, wie wir früher gesehen haben, eine Einheitswurzel ist. Zähler und Nenner einer gebrochenen Invariante müssen daher selbst, wenn auch nur relative, Invarianten sein.

Wenn wir aber nicht gerade die einfachste Darstellung suchen, so können wir absolute gebrochene Invarianten auch als Quotienten von absoluten ganzen Invarianten darstellen. Wir brauchen den Bruch nur durch eine geeignete Potenz des Nenners zu erweitern, also wenn die $\alpha$ $e^{\text{te}}$ Einheitswurzeln sind,
\[
\frac{F(x)}{F_1(x)} = \frac{F(x) F_1(x)^{e-1}}{[F_1(x)]^e}
\]
zu setzen. Hiernach können wir auch alle relativen Invarianten mit einem bestimmten Factorensysteme $\alpha, \beta, \ldots$ darstellen als Product von einer von ihnen mit absoluten Invarianten, die aber gebrochen sein können.

\begin{center}
\S.\,43.

\gesperrt{Das Formenproblem.}
\end{center}

Im vorigen Paragraphen ist nachgewiesen, dass es zu einer endlichen Gruppe linearer Substitutionen eine endliche Anzahl unabhängiger Invarianten giebt. Ist $n$ die Dimension der linearen
%CONT>>
